\documentclass[10pt,a4paper]{amsart}
\usepackage[margin=19mm]{geometry}
\usepackage{amsmath,amssymb,mathtools}
\usepackage[scr=rsfs,cal=euler]{mathalfa}
\usepackage{xfrac}
\usepackage[unicode,hidelinks,bookmarks=false]{hyperref}
\newcommand{\ie}{i.e.\ }
\newcommand{\cf}{cf.\ }
\newcommand{\resp}{respectively\ }
\newcommand{\Subsection}{\subsection}
\providecommand{\nobreakdash}{\nobreak\hbox{-}}
\setlength{\emergencystretch}{2em}
\multlinegap0pt
\usepackage{amssymb}
\usepackage{tikz}
\usepackage{xcolor}
\newtheorem{lemma}{Lemma}
\def\Z{\mathbb{Z}}
\def\codim{\mathrm{codim}}  
\def\um{{\bf m}}
\title{Stable toric sheaves. I \\ Chern classes}
\author{Carl TIPLER}
\date{Source: arXiv:2510.14651v3 (CC BY 4.0)}
\begin{document}
\maketitle

\noindent\emph{Source and licence.} Stable toric sheaves. I \\ Chern classes. Version arXiv:2510.14651v3 (CC BY 4.0), distributed by arXiv under the Creative Commons Attribution 4.0 International licence, http://creativecommons.org/licenses/by/4.0/, as recorded in the arXiv licence metadata for this exact version and shown on the abstract page https://arxiv.org/abs/2510.14651v3. Full licence notice: \texttt{licenses/arxiv-2510.14651-CC-BY-4.0.txt}.\par\smallskip

\noindent\textit{Editorial scope.} Counted unit: the article's lemma cones (source0.tex lines 1543--1550). The article's complete original proof of it is delivered with it (source0.tex lines 1551--1611, one \texttt{proof} environment of 61 lines). No mathematics is written, repaired or re-derived: the statement and the proof are the article's own lines, in the article's own notation, and no retained line is substituted or re-flowed.  No further source-local context is needed: every cross-reference of every retained line resolves inside the two delivered spans.  Nothing was omitted from any retained span, so the package carries no omission record.  No retained line contains a citation, so the wrapper carries no bibliography and no bibliography entry.  No retained line contains a figure environment, an image-inclusion command or drawing code, so no image is delivered and the package declares no asset dependency.  Editorial formatting only, in addition: an AMS \texttt{amsart} preamble that restates the article's own declarations and macros used by the retained lines, namely the environments \texttt{lemma} on separate counters, together with the standard packages those lines need.  The retained environments print in this document as Lemma 1 in order of appearance, whereas the article prints the same items with the numbering of its own full document.  The complete native source of the article as distributed in the arXiv e-print for this exact version is bundled unchanged as \texttt{sources/187199/source0.tex}.
 \begin{lemma}
  \label{lem:sum minus one on cones}
  With the previous notations we have, for $p\leq \codim(\sigma_0)$,
\begin{equation}
 \label{eq:summinusonemsigma}
\sum_{\sigma_0\leq\sigma} (-1)^{\codim(\sigma)}\um_\sigma^p=\um_\Sigma^p.
\end{equation}
 \end{lemma}

 \begin{proof}
 The identity will actually be proved for any set of integers $(m_\rho)_{\rho\in\Sigma(1)}$, setting 
 $\um_\sigma=\sum_{\rho\in\sigma(1)}m_\rho$ for any $\sigma\in\Sigma$, and $\um_\Sigma=\sum_{\rho\in\Sigma(1)} m_\rho$ (forgetting that the $m_\rho$'s come from an elementary injection). 
  We first deal with the case $p=0$. In that case, the proof follows from the following :
  $$
  \sum_{\sigma_0\leq\sigma} (-1)^{\codim(\sigma)}=\sum_{i=k_0}^n\binom{n+1-k_0}{i-k_0}(-1)^{n-i}=1
  $$
  that can be obtained by noting that there are $\binom{n+1-k_0}{i-k_0}$ cones of dimension $i$ that contain $\sigma_0$. 
%   The equality \eqref{eq:summinusonemsigma} is obtained with similar considerations, using that 
%   $$
%   \um_\sigma=\sum_{\rho\in\sigma(1)} m_\rho,
%  $$
%  and counting the number of times that each $m_\rho$ contributes in the left hand side of \eqref{eq:summinusonemsigma}.
 We then proceed to a proof by induction on $p$. Assume first that there is $\rho\in\sigma_0(1)$ with $m_\rho\neq 0$. Then consider
 $$
 \begin{array}{ccc}
 \displaystyle \sum_{\sigma_0\leq\sigma} (-1)^{\codim(\sigma)}\um_\sigma^p & = &\displaystyle\sum_{\sigma_0\leq\sigma} (-1)^{\codim(\sigma)}\Big(\sum_{\rho\in\sigma(1)}m_\rho\Big)^p\\
    & = &\displaystyle\sum_{\sigma_0\leq\sigma} (-1)^{\codim(\sigma)}\Big(\um_{\sigma_0}+\sum_{\rho\in\sigma(1)\setminus\sigma_0(1)}m_\rho\Big)^p\\
     & = &\displaystyle\sum_{\sigma_0\leq\sigma} (-1)^{\codim(\sigma)} \Big(\sum_{q=0}^p \binom{p}{q} \um_{\sigma_0}^{p-q}\Big(\sum_{\rho\in\sigma(1)\setminus\sigma_0(1)}m_\rho\Big)^q \Big)\\
        & = &\displaystyle \sum_{q=0}^p \binom{p}{q}\um_{\sigma_0}^{p-q}\Big(\sum_{\sigma_0\leq\sigma} (-1)^{\codim(\sigma)} \Big(\sum_{\rho\in\sigma(1)\setminus\sigma_0(1)}m_\rho\Big)^q \Big).
 \end{array}
 $$
 Now, for $q<p$, we can use the induction hypothesis on the set of integers $(m_\rho')_{\rho\in\Sigma(1)}$ defined by $m_\rho'=m_\rho$ if $\rho\notin\sigma_0(1)$ and $m_\rho'=0$ if $\rho\in\sigma_0(1)$. This gives
 $$
 \begin{array}{ccc}
 \sum_{\sigma_0\leq\sigma} (-1)^{\codim(\sigma)}\um_\sigma^p & = & \displaystyle \sum_{q=0}^{p-1} \binom{p}{q}\um_{\sigma_0}^{p-q}\Big(\um_\Sigma-\um_{\sigma_0}\Big)^q\\
 & & \displaystyle+\sum_{\sigma_0\leq\sigma} (-1)^{\codim(\sigma)} \Big(\sum_{\rho\in\sigma(1)\setminus\sigma_0(1)}m_\rho\Big)^p .
 \end{array}
 $$
By induction on the number of $\rho$'s such that $m_\rho\neq 0$, we also obtain 
$$
\sum_{\sigma_0\leq\sigma} (-1)^{\codim(\sigma)} \Big(\sum_{\rho\in\sigma(1)\setminus\sigma_0(1)}m_\rho\Big)^p= (\um_\Sigma-\um_{\sigma_0})^p,
$$
and we can conclude in that case that 
$$
\sum_{\sigma_0\leq\sigma} (-1)^{\codim(\sigma)}\um_\sigma^p=\displaystyle \sum_{q=0}^{p} \binom{p}{q}\um_{\sigma_0}^{p-q}(\um_\Sigma-\um_{\sigma_0})^q=(\um_{\sigma_0}+\um_\Sigma-\um_{\sigma_0})^p=\um_\Sigma^p.
$$
What remains is to initialize the last induction used, that is to deal with the case where for each $\rho\in\sigma_0(1)$, $m_\rho=0$. So we are left with showing :
$$
\sum_{k=\dim(\sigma_0)}^n \left(\sum_{\sigma_0\leq\sigma,\,\sigma\in\Sigma(k)} (-1)^{n-k}\Big(\sum_{\rho\in\sigma(1)\setminus\sigma_0(1)}m_\rho\Big)^p\,\right)=\Big(\sum_{\rho\in\Sigma(1)\setminus \sigma_0(1)}m_\rho\Big)^p.
$$
Re-indexing, setting $\codim(\sigma_0)=r$, it is enough to show that for any $(r+1)$-tuple $(m_j)_{0\leq j\leq r}\in\Z^{r+1}$ (that corresponds to $(m_\rho)_{\rho\notin \sigma_0(1)}$) and any $1\leq p \leq r$, we have
\begin{equation}
 \label{eq:alternating sum mis}
\sum_{i=1}^{r} (-1)^i \sum_{\substack{J \subset \{0,\dots,r\} \\ |J|=i}} \Big(\sum_{j \in J} m_j\Big)^p=(-1)^r \Big(\sum_{j=0}^{r} m_j \Big)^p.
\end{equation}
To show this, consider the generating function
$$
F(y) = \sum_{J \subset \{0,\dots,r\}} (-1)^{|J|} e^{y \sum_{j \in J} m_j} = \prod_{j=0}^{r} (1 - e^{y m_j}).
$$
Formally, we can develop in power series :
$$
e^{y \sum_{j \in J} m_j} = \sum_{q=0}^{\infty} \frac{y^q}{q!} \Big(\sum_{j \in J} m_j\Big)^q,
$$
and compute the coefficient in front of $\displaystyle\frac{y^p}{p!}$ in $F(y)$ :
$$
\sum_{J \subset \{0,\dots,r\}} (-1)^{|J|} \Big(\sum_{j \in J} m_j\Big)^p = \sum_{i=0}^{r+1} (-1)^i \sum_{\substack{J \subset \{0,\dots,r\} \\ |J|=i}} \Big(\sum_{j \in J} m_j\Big)^p.
$$
On the other hand, $y$ divides each factor $(1-e^{y m_j})$, hence $F(y)=\prod_{j=0}^{r} (1 - e^{y m_j})$ starts in degree $y^{r+1}$.  
Thus, for $1 \leq p \leq r$, the coefficient in front of $y^p$ in $F(y)$ vanishes, which gives \eqref{eq:alternating sum mis}.
 \end{proof}

\end{document}
